\section{Introduction}\label{sec:intro}

Self-attention \citep{vaswani2017attention} computes an interaction between
every pair of positions in a sequence. That property is the reason
transformers model long-range structure well, and it is also the reason their
cost grows as $\Theta(n^2)$ in the sequence length $n$. For the context
lengths that are now routinely requested of production models --- tens of
thousands of tokens of source code, a whole scientific article and its
appendices, a multi-day conversation --- the attention term stops being a
constant factor and starts being the budget.

The literature has answered this in two broadly different ways. One family
\emph{restricts} the interaction pattern: sliding windows, dilated windows,
strided or block-sparse masks, and hand-designed combinations of the three
\citep{child2019sparse,beltagy2020longformer,zaheer2020bigbird}. These methods
are cheap and predictable, but the pattern is fixed before the model sees any
data, so a dependency the designer did not anticipate simply cannot be
expressed. The other family \emph{approximates} the softmax kernel, replacing
the exact attention matrix with a low-rank or feature-map surrogate
\citep{wang2020linformer,katharopoulos2020linear,choromanski2021performer}.
These methods keep a formally global receptive field, but the approximation
error is not under the practitioner's control and it interacts badly with the
sharply peaked attention distributions that emerge late in training.

We take a third position, which we think is the pragmatic one: keep attention
\emph{exact}, and make the \emph{choice of what to attend to} learned. Each
query is allocated a fixed bandwidth --- a corridor --- of $k$ key positions.
A lightweight scorer, operating on block summaries rather than on individual
positions, decides which blocks the corridor passes through. Because the
bandwidth is a hyperparameter rather than an emergent property, the cost of the
layer is known before training starts, and because the routing is learned, the
pattern can differ per head, per layer, and per input.

Three design commitments follow from this position, and they are what
distinguishes Corridor Attention from the closest prior work.

\paragraph{Bandwidth is declared, not discovered.} The practitioner sets $k$.
The model may not exceed it, and the training objective contains no term that
trades accuracy against cost --- there is nothing to trade, because the cost is
fixed. This makes capacity planning a matter of arithmetic rather than of
empirical scaling curves.

\paragraph{Routing is scored on blocks, decided per query.} Scoring every
candidate key position would reintroduce the quadratic term we are trying to
remove. Corridor Attention scores $\lceil n / b \rceil$ block summaries
instead, at a cost of $\Theta(n^2 / b^2)$ for the scorer, which for the block
sizes we use is a small fraction of the total. The selection itself is
per-query, so two queries in the same block may route differently.

\paragraph{The local window is not negotiable.} Ablating the dense local
window degrades perplexity far more than ablating any equivalent amount of
routed bandwidth. We therefore treat the window as part of the architecture
rather than as one option among several, and we spend the learned bandwidth
only on what the window cannot reach.

Our contributions are as follows. (i) We formalise bounded-bandwidth routed
attention and give a block-scored, per-query instantiation of it that trains
end to end without an auxiliary load-balancing loss. (ii) We give a cost model
for the layer and validate it against measured wall-clock time, showing that
the predicted and observed crossover points against dense attention agree
closely. (iii) We evaluate on long-context language modelling and on synthetic
retrieval probes, and we report the ablations that identify which of the three
components each benefit comes from.

